\documentclass[conference]{IEEEtran}

\usepackage{graphicx}
\usepackage{amsmath,amssymb}
\usepackage{float}
\usepackage{booktabs}
\usepackage{multirow}

\begin{document}

\title{
    \textbf{Benchmark Report: 500.perlbench\_r} \\
    \text{Register Spill Analysis on RISC-V via gem5}
}

\author{
    \IEEEauthorblockN{Lütfullah Burak Kaya}
    \IEEEauthorblockA{
        Ozyegin University \\
        burak.kaya.50711@ozu.edu.tr
    }
    \and
    \IEEEauthorblockN{Ismail Akturk}
    \IEEEauthorblockA{
        Ozyegin University \\
        ismail.akturk@ozyegin.edu.tr
    }
}

\newcommand{\LongDate}{%
    \number\day\space
    \ifcase\month
    \or January\or February\or March\or April\or May\or June%
    \or July\or August\or September\or October\or November\or December%
    \fi
    \space \number\year
}

\twocolumn[
\begin{@twocolumnfalse}
\maketitle
\begin{center}{\small \LongDate}\end{center}
\vspace{1em}
\end{@twocolumnfalse}
]

% ================================================

\begin{abstract}
This report presents the register spill characterization of the
\texttt{500.perlbench\_r} benchmark from SPEC CPU2017, executed on a
RISC-V (RV64GC) target using the gem5 TimingSimpleCPU simulator.
A custom SpillDetector was integrated into gem5 to count store-then-load
patterns on the stack within a user-defined Region of Interest (ROI).
The \texttt{checkspam.pl} workload (2500 messages, single-process) was
used as the input, selected for compatibility with gem5's SE mode which
does not support \texttt{fork()}. Results show a spill rate of
approximately 7.28\% within the ROI --- the highest observed across all
SPEC CPU2017 benchmarks evaluated in this study --- corresponding to over
1 million detected spill events across 14.7 million ROI instructions.
\end{abstract}

% =========================================================
\section{Benchmark Overview}

\texttt{500.perlbench\_r} is the SPEC CPU2017 rate variant of a
stripped-down Perl interpreter. The benchmark compiles and executes
real-world Perl scripts, exercising the interpreter's parser, bytecode
compiler, and runtime engine. It is one of the most complex integer
workloads in the SPEC suite, with a large and irregular code footprint
spanning the Perl core and several extension modules.

The RISC-V binary (\texttt{perlbench\_r.riscv}) was compiled with gem5
ROI markers inserted around the \texttt{perl\_run(my\_perl)} call in
\texttt{src/perlmain.c}. This call is the top-level dispatch into the
Perl interpreter's execution loop and encompasses all script execution:

\begin{itemize}
    \item \texttt{m5\_work\_begin(0,0)} --- before \texttt{perl\_run()}
    \item \texttt{m5\_work\_end(0,0)}   --- after  \texttt{perl\_run()}
\end{itemize}

\subsection{Input Selection}

SPEC CPU2017 provides three dataset sizes for \texttt{perlbench\_r},
each running a different Perl workload (Table~\ref{tab:inputs}).
The standard test dataset (\texttt{test.pl}) runs the Perl internal
test suite by forking one sub-process per test file, which is
incompatible with gem5's SE mode (no \texttt{fork()} support).
Therefore, \texttt{checkspam.pl} from the refrate dataset was used.
This script is a single-process spam filter based on SpamAssassin
rules and exercises the full Perl interpreter without any
process spawning.

\begin{table}[H]
\centering
\caption{perlbench\_r Input Dataset Sizes}
\label{tab:inputs}
\begin{tabular}{lll}
\toprule
\textbf{Dataset} & \textbf{Script} & \textbf{Description} \\
\midrule
test    & test.pl      & Perl internal test suite (fork-based) \\
train   & diffmail.pl  & Email diff utility \\
refrate & checkspam.pl & SpamAssassin-based spam filter \\
\bottomrule
\end{tabular}
\end{table}

The \textbf{checkspam.pl} workload was invoked with the refrate
argument set: \texttt{2500 5 25 11 150 1 1 1 1}, which classifies
2500 email messages against a set of SpamAssassin rules.

% =========================================================
\section{Simulation Setup}

\subsection{gem5 Configuration}

\begin{itemize}
    \item \textbf{CPU model:}   TimingSimpleCPU
    \item \textbf{ISA:}         RISC-V RV64GC
    \item \textbf{Caches:}      enabled (default L1 I/D + L2)
    \item \textbf{Memory:}      4\,GB simulated physical RAM
    \item \textbf{Spill mode:}  summary only (\texttt{spill\_verbose=False})
\end{itemize}

The full gem5 invocation was:

\begin{verbatim}
./build/RISCV/gem5.opt
  --outdir=benchmarks/perlbench/m5out_test
  configs/deprecated/example/se.py
  --cmd=benchmarks/perlbench/perlbench_r.riscv
  --options="-I./benchmarks/perlbench
             -I./benchmarks/perlbench/lib
             benchmarks/perlbench/checkspam.pl
             2500 5 25 11 150 1 1 1 1"
  --cpu-type=TimingSimpleCPU
  --caches
  --mem-size=4GB
\end{verbatim}

% =========================================================
\section{Simulation Results}

The simulation completed with exit code 0. The \texttt{checkspam.pl}
script processed all 2500 email messages and printed per-message
classification results.

\subsection{Pre-ROI Context}

Table~\ref{tab:preroi} shows the global counters at the moment
\texttt{m5\_work\_begin} was triggered.

\begin{table}[H]
\centering
\caption{Global Counters at ROI Entry}
\label{tab:preroi}
\begin{tabular}{lr}
\toprule
\textbf{Metric} & \textbf{Value} \\
\midrule
Total instructions (pre-ROI) & 1,134,635 \\
Total spills detected (pre-ROI) & 65,714 \\
Pre-ROI spill rate & 5.79\% \\
\bottomrule
\end{tabular}
\end{table}

\subsection{ROI Spill Statistics}

Table~\ref{tab:roi} presents the spill statistics collected within
the ROI.

\begin{table}[H]
\centering
\caption{ROI Spill Statistics --- 500.perlbench\_r (checkspam.pl, 2500 msgs)}
\label{tab:roi}
\begin{tabular}{lr}
\toprule
\textbf{Metric} & \textbf{Value} \\
\midrule
ROI instructions         &  14,690,864 \\
ROI stores               &   2,035,472 \\
ROI loads                &   3,491,473 \\
\midrule
\textbf{ROI spills}      & \textbf{1,070,055} \\
\textbf{Spill rate}      & \textbf{7.2838\%} \\
\bottomrule
\end{tabular}
\end{table}

% =========================================================
\section{Analysis}

\subsection{Spill Rate Interpretation}

The measured spill rate of \textbf{7.28\%} means that approximately
1 in every 14 instructions inside the ROI is a register spill.
This is the highest spill rate observed among all SPEC CPU2017
benchmarks evaluated in this study (Table~\ref{tab:compare}).

\begin{table}[H]
\centering
\caption{Cross-Benchmark Spill Rate Comparison}
\label{tab:compare}
\begin{tabular}{lrr}
\toprule
\textbf{Benchmark} & \textbf{ROI Spills} & \textbf{Spill Rate} \\
\midrule
500.perlbench\_r & 1,070,055   & \textbf{7.2838\%} \\
526.blender\_r   & 61,845,869  & 6.2193\% \\
544.nab\_r       & 151,310,616 & 2.2612\% \\
505.mcf\_r       & 446,113,136 & 1.8263\% \\
508.namd\_r      & 327,718,955 & 1.0269\% \\
\bottomrule
\end{tabular}
\end{table}

\subsection{Store vs.\ Load Balance}

The load count (3.49M) is approximately $1.72\times$ higher than the
store count (2.04M). This near-symmetric ratio is characteristic of
interpreter dispatch loops: the interpreter must load operands,
dispatch through a function pointer table (loading the handler
address), and frequently save/restore interpreter state registers
across recursive calls into user-defined Perl functions.

\subsection{Spill Fraction of Loads}

Within the ROI, the ratio of spills to total loads is:
\[
\frac{1{,}070{,}055}{3{,}491{,}473} \approx 30.6\%
\]
Nearly 1 in 3 load instructions is a spill reload. This is the
highest such fraction across all benchmarks evaluated, reflecting the
deep and irregular call stack of the Perl interpreter: each level of
Perl function call pushes interpreter state onto the C stack, creating
severe register pressure in the RISC-V backend.

\subsection{Pre-ROI Spill Density}

Pre-ROI spills: $65{,}714$ over $1{,}134{,}635$ instructions
= \textbf{5.79\%}. This is high relative to most other benchmarks'
pre-ROI phases, indicating that even Perl interpreter startup ---
module loading, script parsing, and bytecode compilation via
\texttt{perl\_parse()} --- is register-pressure-intensive on RISC-V.

\subsection{Interpreter-Driven Spill Pressure}

The extreme spill rate in perlbench is structurally driven by the
Perl interpreter's C implementation. The bytecode dispatch loop
must maintain a large set of live variables simultaneously:
the interpreter's program counter (\texttt{PL\_op}), the stack
pointer (\texttt{PL\_stack\_sp}), the current pad (lexical variable
frame), the current context stack, and various per-opcode temporaries.
With RISC-V providing 32 general-purpose registers (of which only
$\sim$27 are caller/callee-save usable), the compiler cannot keep
all of these live across the frequent \texttt{call}/\texttt{ret} pairs
of opcode handler dispatch, resulting in pervasive register spilling.

% =========================================================
\section{Conclusion}

The \texttt{500.perlbench\_r} benchmark exhibits the highest register
spill rate (7.28\%) among all SPEC CPU2017 benchmarks evaluated in
this study. With 1.07 million spill events across 14.7 million ROI
instructions, approximately 1 in 3 load instructions is a spill reload.
This is attributable to the Perl interpreter's deeply recursive C call
structure, which continuously exceeds the RISC-V register file capacity
during opcode dispatch. These results identify
\texttt{500.perlbench\_r} as the most register-pressure-intensive
benchmark in the evaluated suite and a prime candidate for studying
compiler spill-reduction techniques on RISC-V.

\end{document}
